\documentclass[10pt,a4paper]{amsart}
\usepackage[margin=19mm]{geometry}
\usepackage{amsmath,amssymb,mathtools}
\usepackage[scr=rsfs,cal=euler]{mathalfa}
\usepackage{xfrac}
\usepackage[unicode,hidelinks,bookmarks=false]{hyperref}
\newcommand{\ie}{i.e.\ }
\newcommand{\cf}{cf.\ }
\newcommand{\resp}{respectively\ }
\newcommand{\Subsection}{\subsection}
\providecommand{\nobreakdash}{\nobreak\hbox{-}}
\setlength{\emergencystretch}{2em}
\multlinegap0pt
\usepackage{mathtools}
\usepackage{amssymb}
\usepackage{enumerate}
\usepackage{tikz}
\newtheorem{theorem}{Theorem}
\title{This is the title}
\author{K. Mahesh Krishna}
\date{Source: arXiv:2203.06916v1 (CC BY 4.0)}
\begin{document}
\maketitle

\noindent\emph{Source and licence.} This is the title. Version arXiv:2203.06916v1 (CC BY 4.0), distributed by arXiv under the Creative Commons Attribution 4.0 International licence, http://creativecommons.org/licenses/by/4.0/, as recorded in the arXiv licence metadata for this exact version and shown on the abstract page https://arxiv.org/abs/2203.06916v1. Full licence notice: \texttt{licenses/arxiv-2203.06916-CC-BY-4.0.txt}.\par\smallskip

\noindent\textit{Editorial scope.} Counted unit: the article's theorem CONVEX (source0.tex lines 124--133). The article's complete original proof of it is delivered with it (source0.tex lines 134--155, one \texttt{proof} environment of 22 lines). No mathematics is written, repaired or re-derived: the statement and the proof are the article's own lines, in the article's own notation, and no retained line is substituted or re-flowed.  No further source-local context is needed: every cross-reference of every retained line resolves inside the two delivered spans.  Nothing was omitted from any retained span, so the package carries no omission record.  No retained line contains a citation, so the wrapper carries no bibliography and no bibliography entry.  No retained line contains a figure environment, an image-inclusion command or drawing code, so no image is delivered and the package declares no asset dependency.  Editorial formatting only, in addition: an AMS \texttt{amsart} preamble that restates the article's own declarations and macros used by the retained lines, namely the environments \texttt{theorem} on separate counters, together with the standard packages those lines need.  The retained environments print in this document as Theorem 1 in order of appearance, whereas the article prints the same items with the numbering of its own full document.  The complete native source of the article as distributed in the arXiv e-print for this exact version is bundled unchanged as \texttt{sources/123190/source0.tex}.
\begin{theorem}
\textbf{(C*-algebraic Gauss-Lucas Theorem)} Let $p(z)=(z-a_1)(z-a_2)\cdots (z-a_n)\in\mathcal{A}[z]$. 	Define 
\begin{align*}
	\mathcal{B}\coloneqq \{z \in \mathcal{A}: z-a_j \text{ is  invertible in }  \mathcal{A} \text{ for all } 1\leq j \leq n\}, \quad 	\mathcal{C}\coloneqq \{z \in \mathcal{A}: p'(z)=0\}.	
\end{align*}
 Then for each $z \in \mathcal{B}\cap  \mathcal{C}$, there are positive $\omega_{z_1}, \dots, \omega_{z_n} \in  \mathcal{A}$ such that 
\begin{align}\label{CONVEX}
	z=\sum_{j=1}^{n}\omega_{z_j}a_j, \quad \sum_{j=1}^{n}\omega_{z_j}=1.
\end{align}
\end{theorem}

\begin{proof}
	We have 
	\begin{align*}
		p'(z)=\sum_{j=1}^{n}(z-a_1)\cdots \widehat{(z-a_j)}\cdots (z-a_n), \quad \forall z \in \mathcal{A}.
	\end{align*}
	where the term with cap is missing. Then 
	\begin{align*}
		p(z)^{-1}	p'(z)=\sum_{j=1}^{n}(z-a_j)^{-1}, \quad \forall z \in \mathcal{B}.
	\end{align*}
We therefore have 
\begin{align}\label{FIRST}
	0=\sum_{j=1}^{n}(z-a_j)^{-1}=\sum_{j=1}^{n}((z-a_j)(z-a_j)^*)^{-1}(z^*-a_j^*), \quad \forall z \in \mathcal{B}\cap  \mathcal{C}.
\end{align}
Now  define 
\begin{align*}
	\omega_z \coloneqq \sum_{j=1}^{n}((z-a_j)(z-a_j)^*)^{-1}, \quad \forall z \in \mathcal{B}\cap  \mathcal{C}.
\end{align*}
Note that whenever $ z \in \mathcal{B}$, each $((z-a_j)(z-a_j)^*)^{-1}$ is positive and invertible, $1\leq j \leq n$. Hence $\omega_z$ is positive and invertible. Equation (\ref{FIRST}) then gives 
\begin{align*}
z=\omega_z^{-1}\sum_{j=1}^{n}((z-a_j)(z-a_j)^*)^{-1}a_j, \quad \forall z \in \mathcal{B}\cap  \mathcal{C}.
\end{align*}
\end{proof}

\end{document}
